% !TeX root = ../papers/nc-p3.tex
\section{有限四辺と四次中心列を持つ窓の二つの中心次数}
\label{sec:central-cut-center-dichotomy}
\begin{theorem}\label{thm:central-cut-center-dichotomy}
標準正則な有限四辺の窓で、Frobenius標準折り畳み $B$ が四つの中心四次パラメータの
多項式環 $P$ 上有限自由とする。この仮定は既出の全次数中心切断検査で認証できる。
全中心が素、正規またはGorensteinであることを追加する必要はない。
$B$ は素、全中心 $Z$ は整かつ正規なGorenstein環で有理特異点を持ち、
$B$ は中心の正則部分でAzumayaである。換算次数で $\omega_Z=Z(-1)$ である。
$d=\operatorname{rank}_P Z$、$s$ を頂点ごとの可除指数とすると、可能性は
\[
\begin{array}{c|c|c|c}
 d&s&\dim Z_1& H_{Z/P_+Z}(t)\\\hline
 4&4&5&1+t+t^2+t^3\\
 16&2&11&1+7t+7t^2+t^3 .
\end{array}
\]
$d=4$ なら全中心は普通の四次超曲面、または整な三次元二次曲面の四次分岐二重被覆である。
二次曲面の滑らかさは仮定しない。$d=16$ なら一般超平面は滑らかなK3、
$d=4$ なら一般超平面は滑らかなK3または有限個の $A_1$ を持つK3であり、
各 $A_1$ に現れる頂点表現は二つの指標を二コピーずつ持つ。
\end{theorem}
\begin{proof}[素性と相対双対]
頂点射影 $e_iB$ の非零な次数付き部分加群は、$P$ 上torsion-freeなので正の $P$ 階数を持つ。
商も正の $P$ 階数を持つことはない。実際、有限生成次数付き加群のHilbert級数は
有限射影分解から $p(t)/(1-t)^4$ であり、$P$ 上正階数なら $p(1)$ は正整数である。
二つの $p(1)$ の和が頂点射影の1になることに反する。非零の斉次経路元の左乗法は
頂点射影間の写像で、両端は $P$ 上同階数、像は全階数なので単射である。
これから零でない合成可能な斉次元の積は非零で、最高次を取れば非斉次の各隅でも同じである。
任意の二つの非零元では冪等元で非零隅を取り、その間には非零な経路があるため $B$ は素である。
有限生成の中心上有限なので一般商は中心単純である。

$D=\operatorname{Hom}_P(B,P)$ とする。中心正則列を割った有限次元代数の
Frobenius性から、$D$ は左加群として $B(12)$ に同型である（ここだけは経路長次数）。
実際商の最高次12の汎関数を斉次自由基底で持ち上げ、積対合が剰余で完全であることを使う。
右でも同じ汎関数が同型を与えるので、相対Nakayama自己同型を持つ可逆両側加群である。
全ての閉点単純に対し $\operatorname{Ext}^j_B(L,D)=\operatorname{Ext}^j_P(L,P)$ の非零次数は4である。
大域次元は高々4であるから $B$ はhomologically homogeneousで、
Stafford--Van den Berghの定理により中心は正規で有理特異点を持つ。

さらに標準折り畳みで $D\simeq B(12)$ は\emph{両側}同型である。詳細を与える。
$E=B^!$ の斉次四次traceは標準周期の定義により通常の積について対称である
（または命題\ref{prop:general-cut-trivial-extension}の自明拡大traceを取る）。
$W_i=(E_i)^*$ とするKoszul両側分解の最終双対余核は
$\operatorname{Ext}_{B^e}^4(B,B^e)=B(4)$ である。ここで $B^e=B\otimes_k B^{\rm op}$。
実際 $W_4$ の各生成元の双対を始点ごとに $g_i$ とすると、
$W_3$ は反対の一文字にtraceで双対である。最終微分は左端の移動と
$(-1)^4$ 倍の右端の移動の和なので、双対余核の全関係は
一文字が $g$ を通ると符号を変える関係である。両係数が同じなのは対称traceによる。
$g_i$ に頂点の交代符号を掛ければ各一文字と可換になり、これが正則両側加群の表示になる。

$P^e=P\otimes_k P$ とし、$B^e$ が $P^e$ 自由であることから
\[
 \operatorname{Ext}^4_{B^e}(B,\operatorname{Hom}_{P^e}(B^e,P^e))
 =\operatorname{Ext}^4_{P^e}(B,P^e)=D(16)
\]
を得る。右側は次数4の四つの変数差で対角 $P$ を解消した後、$B$ の $P$ 自由性を使ったもの。
左の第二引数は外側と内側の作用を備えた二つの $D$ の外積である。
有限両側射影分解の双対にこの外積を掛けて計算すると、
内側作用を残した結果は
$D\otimes_B\operatorname{Ext}^4_{B^e}(B,B^e)\otimes_B D$ となる。
各有限射影項でのテンソル・Hom同型と $D$ の両側可逆性による。
よって $D\otimes_B D(4)=D(16)$、可逆因子を打ち消して主張が従う。
\end{proof}

\begin{proof}[中心と一般切断]
被約traceは正規性により $B\to Z$ に値を持ち、中心への制限は非零整数倍なので分裂全射である。
従って $Z$ は $P$ 自由加群の次数付き直和因子で、次数付き自由である。
$D$ 内の中心部分（巡回な $P$ 線型汎関数）は $\operatorname{Hom}_P(Z,P)$ に一致する。
一般分数体では各巡回汎関数は被約traceとの合成であり、$Z$ への制限を行列次数で
割ればその因子は整係数に戻る。逆にこの合成は巡回である。
よって前段の両側同型の中心を取ると $\operatorname{Hom}_P(Z,P)=Z(12)$ となる。
$\omega_P=P(-16)$ を掛ければ標準加群は $Z(-4)$、換算次数では $Z(-1)$ である。
さらに $D$ の両側生成元である巡回汎関数がこの中心双対の生成元と被約traceの合成なので、
$B$ の被約trace対合は $Z$ 上完全である（$P$ への随伴で同型を検出できる）。
命題\ref{prop:smooth-azumaya}の分離元の証明により、中心正則部分はAzumayaである。

$Z_0=k$ は全矢が非零で四頂点が連結であることから従う。
$Z/P_+Z$ は従ってsocle次数3の可換斉次Frobenius環で、四係数は $(1,b,b,1)$ である。
前定理の候補で $d=1$ はこれに反し、$d=64$ は有限辺に反する。
残りが表の二行である。$d=4$ で剰余の一次生成元の平方が非零なら
完全対合によりその三乗も非零で、四次モニック一元表示になる。
平方が零なら二次生成元を一つ補い、四基底 $1,z,y,zy$ の積から
$z^2=a_1z+a_2$, $y^2=b_1zy+b_2y+b_3z+b_4$ を得る。
二式ともそれぞれでモニックであり階数比較で全表示である。
特に $y$ を平行移動すれば整な二次曲面上の四次分岐二重被覆である。

最後に一般超平面は基点無し豊富な $H$ に属し、幾何学的に整かつ正規で、
双対化層は随伴で自明、$H^1(\mathcal O)=0$ は自由な中心環の切断列から従う。
定理\ref{thm:central-cut-transverse-a-one}の普遍接続・有理二重点・反射的生成元の証明は
次数や二次曲面の滑らかさを使わず、この $H$ についても適用できる。
$d=16$ なら $s=2$ であり、特異点があると $2=\sum d_j\dim\rho_j$、
全て $d_j\ge2$ で少なくとも二種類の表現が必要となり矛盾する。
$d=4$ は同定理の通りである。特にこの議論は、中心切断に通った有限窓について個々に適用できる。
\end{proof}
